\documentclass[11pt]{ctexart}
\usepackage[a4paper,margin=1in]{geometry}
\usepackage{amsmath,amssymb,amsthm,mathtools,bm}
\usepackage{hyperref}
\usepackage{booktabs}
\usepackage{enumitem}

\hypersetup{
  colorlinks=true,
  linkcolor=blue,
  urlcolor=blue,
  citecolor=blue
}

\title{TurboQuant：具有近最优失真率的在线向量量化}
\author{
  Amir Zandieh \and
  Majid Daliri \and
  Majid Hadian \and
  Vahab Mirrokni
}
\date{}

\begin{document}
\maketitle

\begin{abstract}
向量量化（Vector Quantization, VQ）起源于香农的有损源编码理论，目标是在压缩高维欧氏向量时尽量减少几何结构失真。本文提出 \textbf{TurboQuant}，同时优化均方误差（MSE）与内积失真，解决现有方法难以达到最优失真率的问题。

本文算法是\textbf{数据无关（data-oblivious）}且适合在线场景：在任意比特宽度与维度下，都能达到接近最优（仅差一个小常数因子）的失真率。核心思想是：先对输入向量做随机旋转，使坐标服从集中化的 Beta 分布；再利用高维下不同坐标近似独立的性质，对各坐标分别使用最优标量量化器。

针对内积估计，作者指出仅以 MSE 最优为目标会引入偏差。因此提出两阶段方案：先做 MSE 量化，再对残差施加 1-bit QJL（Quantized Johnson-Lindenstrauss）变换，从而得到无偏内积量化器。

此外，论文给出信息论下界证明，说明任意向量量化器的最佳失真率存在极限，TurboQuant 与该下界仅差约 2.7 倍常数。实验显示：在 KV Cache 量化中，3.5 bit/通道可达到几乎无损质量，2.5 bit/通道仅有轻微退化；在近邻搜索任务中，方法在召回率上优于现有 PQ，并将索引时间降至近乎零。
\end{abstract}

\section{引言}
向量量化在高维计算中至关重要，广泛用于大模型推理与向量检索系统。核心任务是把浮点向量压缩为低比特整数表示，同时尽量保持几何信息（如距离与内积）。

在大语言模型部署中，推理延迟常受 HBM/SRAM 间通信瓶颈制约。对权重和激活做量化可显著降低推理开销。尤其在 Transformer 解码阶段，KV Cache 随模型规模和上下文长度增长很快，成为显存与速度瓶颈，因此需要在不明显损伤精度的前提下压缩 KV。

向量数据库中的近邻搜索同样依赖高效量化。现有方法通常二选一：要么计算慢、不利于加速器并行；要么失真率相对于比特宽度不够理想。本文目标即设计一个轻量、在线可用、对 GPU/TPU 友好的量化算法。

TurboQuant 的核心为两阶段：先实现 MSE 近最优的向量量化，再对残差做 1-bit 量化，实现内积无偏且低失真估计。

\section{问题定义}
量化映射定义为
\[
Q: \mathbb{R}^d \to \{0,1\}^B,\quad B=b\cdot d,
\]
其中 $b$ 是每个维度平均比特数（bit-width）。解码映射为
\[
Q^{-1}: \{0,1\}^B \to \mathbb{R}^d.
\]

关注两类失真：
\begin{align}
D_{\text{mse}} &= \mathbb{E}\!\left[\|x-Q^{-1}(Q(x))\|_2^2\right], \\
D_{\text{prod}} &= \mathbb{E}\!\left[\left|\langle y,x\rangle-\langle y,Q^{-1}(Q(x))\rangle\right|^2\right].
\end{align}
并要求内积估计满足无偏性：
\[
\mathbb{E}\!\left[\langle y,Q^{-1}(Q(x))\rangle\right]=\langle y,x\rangle.
\]
论文在最坏情形输入向量下分析，量化器允许随机化。

\section{相关工作}
\begin{itemize}[leftmargin=1.6em]
  \item \textbf{经典 VQ 理论}：从 Shannon、Zador、Gersho 的失真率理论发展而来，但早期算法实现成本高。
  \item \textbf{在线 vs 离线量化}：离线方法需大量数据依赖校准，不适合动态在线场景；在线方法无需数据拟合，更适合实时系统。
  \item \textbf{KV Cache 压缩}：包括结构改造、token 裁剪/驱逐、以及量化方法；QJL 提供了数据无关的一比特内积无偏估计基础。
  \item \textbf{近邻检索中的 PQ}：很多 PQ 依赖 k-means 训练码本，不适合在线；近期无训练方案理论与工程上仍有不足。
\end{itemize}

\section{技术路线与主要贡献}
\subsection{MSE 优化版 TurboQuant}
先对输入向量做随机旋转。旋转后每个坐标服从 Beta 分布，高维下近似高斯，且不同坐标近似独立。因此可以把高维向量量化近似分解为逐坐标标量量化问题。

作者通过求解一维连续 k-means（Lloyd-Max）得到最优标量量化码本，并预先离线计算各常用 bit-width 的码本供在线调用。理论结果（单位范数向量）：
\[
D_{\text{mse}}(Q_{\text{mse}})\le \frac{\sqrt{3}\pi}{2}\cdot 4^{-b}.
\]
在 $b=1,2,3,4$ 时，给出更精细估计：约 $0.36, 0.117, 0.03, 0.009$。

\subsection{内积优化版 TurboQuant}
仅 MSE 最优会导致内积有偏。为此提出两阶段方案：
\begin{enumerate}[leftmargin=1.6em]
  \item 先用 $b-1$ bit 的 $Q_{\text{mse}}$ 量化；
  \item 对残差用 1-bit QJL 量化。
\end{enumerate}
最终达到内积无偏，并获得低失真上界：
\[
D_{\text{prod}}(Q_{\text{prod}})
\le
\frac{\sqrt{3}\pi^2\|y\|_2^2}{d}\cdot 4^{-b}.
\]
在 $b=1,2,3,4$ 时，细化为
\[
\frac{1.57}{d},\ \frac{0.56}{d},\ \frac{0.18}{d},\ \frac{0.047}{d}.
\]

\subsection{信息论下界}
结合 Shannon 下界与 Yao 极小极大原理，论文证明任意随机量化器都存在困难输入使得
\[
D_{\text{mse}}\ge 4^{-b},\qquad
D_{\text{prod}}\ge \frac{\|y\|_2^2}{d}\cdot 4^{-b}.
\]
因此 TurboQuant 与理论最优仅差常数因子（约 2.7），在低比特下更接近最优（如 $b=1$ 约 1.45 倍）。

\section{预备知识}
论文给出关键概率事实：单位超球面上均匀随机点的单坐标服从 Beta 分布，高维下逼近 $\mathcal{N}(0,1/d)$。该性质支撑旋转后逐坐标独立量化的合理性。

同时回顾 Shannon Lower Bound（SLB）并给出其在超球面均匀分布情形下的推论，为后续下界证明服务。

\section{TurboQuant 算法细节}
\subsection{MSE 版本}
\textbf{量化（Quant）}：先用随机旋转矩阵 $\Pi$ 变换 $x\to y=\Pi x$，再对每个坐标 $y_j$ 找最近码字索引并存储。

\textbf{反量化（DeQuant）}：按索引恢复各坐标码字，再乘 $\Pi^\top$ 旋回原坐标系得到重建向量。

作者还讨论了对码字索引做熵编码可再降约 5\% 平均比特宽度（尤其在 4 bit），但为保持系统简洁与速度，主方案未启用。

\subsection{内积版本}
在 MSE 版本基础上，先得到重建向量并计算残差 $r=x-\tilde{x}_{\text{mse}}$，再对 $r$ 做 QJL 符号量化并保存符号与残差范数，解码时把 MSE 重建项与 QJL 重建项相加。该设计理论上保证内积无偏，同时利用残差小范数把误差压低。

\section{KV Cache 量化}
论文将方法用于自回归推理时的 KV Cache 压缩：以量化后的 $\hat{K},\hat{V}$ 近似注意力计算。文中示例（$d=128$, fp16）显示 3-bit 量化可带来约 4 倍以上存储压缩。

\section{实验结果}
\subsection{理论验证}
在 DBpedia 向量数据（1536 维）上验证：随着 bit-width 增加，两种方法误差方差都下降；$Q_{\text{prod}}$ 在所有位宽下保持内积无偏；$Q_{\text{mse}}$ 在低位宽时存在偏差，但位宽增大后偏差趋近于 0；实测曲线与理论上/下界吻合良好。

\subsection{Needle-in-a-Haystack 长上下文检索}
在 Llama-3.1-8B-Instruct 上，文档长度覆盖 4k--104k token，压缩比设为 0.25。TurboQuant 在 4$\times$ 压缩下达到与全精度几乎一致（文中给出同分）的召回表现，优于多种对比方法。

\subsection{LongBench 端到端生成}
在 Llama 与 Ministral 上测试 2.5bit / 3.5bit。3.5bit 结果与全精度平均分基本持平，2.5bit 仅轻微退化；并且方法在生成流过程中也执行量化，仍保持竞争力。

\subsection{近邻搜索}
与 Product Quantization、RabitQ 对比时，TurboQuant 在召回上更优，同时量化（建索引）时间极低，体现了在线/向量化实现优势。

\section{结论}
TurboQuant 提供了一个兼顾理论与工程可用性的在线向量量化框架：对 MSE 与内积估计均有近最优失真率；内积版本可做到无偏；常数因子接近信息论下界；在 KV Cache 与向量检索中展现出高压缩、低开销和高精度。

从实用角度看，TurboQuant 的价值在于：无需数据依赖训练码本、可在线部署、天然适合加速器并行，尤其适合长上下文推理与大规模向量检索系统。

\end{document}
